\documentclass[10pt]{article}
\usepackage[T1]{fontenc}
\usepackage{lmodern}
\usepackage[margin=2.2cm]{geometry}
\usepackage{booktabs,tabularx,array}
\usepackage{xcolor}
\usepackage[hidelinks]{hyperref}
\setlength{\parindent}{0pt}
\setlength{\parskip}{5pt}
\newcommand{\rc}[1]{\par\textit{Comment: #1}\par}
\newcommand{\resp}{\textbf{Response.} }

\begin{document}

\begin{center}
{\Large Response to the Reviewers}\\[4pt]
Paper \#1571347280, ICDABI 2026 (DATA'26)\\
Nahla Nabil Skaik, Faculty of Computer Studies, Arab Open University, Kingdom of Bahrain
\end{center}

I thank the three reviewers for their careful reading and detailed comments. The revision addresses every point raised, and in doing so it changed the paper substantially. The main changes are summarized first, followed by point-by-point responses. Section, table and figure numbers refer to the revised manuscript.

\textbf{Revised title.} \emph{Physics-Informed Graph Neural Networks with Uncertainty Quantification for Conservative, Registry-Compatible Blue Carbon Crediting.} The earlier subtitle presented the evaluation protocol as the contribution; following Reviewer~1, the protocol is now the method of evidence and the contribution is the crediting pipeline and its empirical tests.

\textbf{Main changes.}
\begin{enumerate}\setlength{\itemsep}{1pt}
\item \textbf{A data error was found and corrected.} While addressing the negative $R^2$, I found that the Sentinel-2 clip used in the submitted version was shifted 9.18\,km north of its true position (the mean NDVI of mangrove nodes was 0.003). Re-applying the same shift to correctly referenced imagery reproduced the submitted features (correlation 0.91--0.97). All features were rebuilt from correctly referenced imagery (mean NDVI 0.46). With every other setting unchanged, this correction alone raised the GNN's cross-validated $R^2$ from 0.17 to 0.26 (Section~III-A, Section~IV-A).
\item \textbf{Stronger and fairer models:} random forest, XGBoost, spatial-lag and location-only baselines, and a physics-informed GNN with a new hydroperiod prior. The selected ensemble of the physics-informed GNN and XGBoost reached test $R^2=0.28$ and 0.60 in the two regions (Table~I).
\item \textbf{A second region and a second date}: the Everglades coast (22,326 nodes) and August 2026 imagery (Section~III-A, IV-B).
\item \textbf{A registry rule instead of an ad hoc multiplier}: the Verra VM0033 uncertainty deduction (Eq.~2), kept separate from the non-permanence buffer and checked against registry records and observed mangrove loss (Section~III-C, IV-C).
\item \textbf{Independent validation}: 87 field soil cores from the Smithsonian Coastal Carbon Library and a field-calibrated regional map (Section~IV-D).
\item \textbf{Presentation}: a pipeline figure (Fig.~1) now shows how calibrated intervals become the registry deduction, and the business effect is quantified (Section~IV-C).
\item \textbf{Protocol}: all selection rules were committed before the experiments they govern. The test split was read once for all pre-specified models; one further model (the ensemble), chosen on cross-validation alone, was evaluated in a second, declared read. Both are reported (Section~III-D, Table~I). Code, rules and every result file are public: \url{https://github.com/Nahla-Nabil/PIGNN-UQ}.
\end{enumerate}

\section*{Reviewer 1}

\rc{The substantive contribution is thin \dots\ the claimed main contribution, a validation/test protocol, is ordinary good practice rather than a methodological advance. Each component is standard.}
\resp I agree that the components are standard and that a protocol alone is not a contribution. The revision reframes the paper around four contributions (Section~I): (1) per-node calibrated intervals converted into the VM0033 uncertainty deduction and compared with buffers held by issuing projects; (2) an empirical test of the methodology's aggregation rule, which covered only 18\% and 33\% of realized project-level errors because spatial model errors are correlated (Table~II); (3) validation against 87 field cores; and (4) two regions and two dates with leakage-free splits. A research-gap paragraph now closes Section~II.

\rc{All models have negative test $R^2$, so they do worse than predicting the mean.}
\resp The negative values were caused largely by the georeferencing error described above. After correction, test $R^2$ is 0.28 (Region~1) and 0.60 (Region~2) for the selected ensemble, 0.20 and 0.59 for XGBoost, against $-0.11$ and $-0.01$ for the training mean (Table~I). Results on interior test nodes, two hops from any other split, are also reported (0.19 and 0.58).

\rc{The ``prior-generalization'' advantage is largely circular.}
\resp Agreed. The prior is no longer used as evidence. It was replaced by a physically grounded hydroperiod prior (tidal inundation frequency from lidar elevation and a year of tide-gauge records, plus a red-edge index), and its correlation with the predictions it constrains is now 0.14 (against 0.76 before), so the model no longer copies it. Evidence comes instead from independent field cores (Section~IV-D).

\rc{The CQR intervals, which do not even use the physics-informed model, are several times wider than the label's standard deviation.}
\resp A quantile version of the physics-informed GNN was added, and crediting is reported under all five interval methods (credited stock 0.84--0.90 and 0.95--0.98 of the buffer-only amount). On width: a 90\% interval of a variable with standard deviation $\sigma$ spans about $3.29\sigma$, so the right reference is a trivial conformal interval from a constant predictor. Against that reference, the selected intervals are 32\% narrower in Region~2 and as wide in Region~1, where they score 18\% better on the Winkler score and their width ranks the actual errors (Spearman 0.54). Section~V shows that the achievable narrowing is bounded by accuracy (about 35\% in Region~2, where 32\% was reached).

\rc{The pricing mechanism rests on arbitrary parameters and conflates estimation uncertainty with non-permanence risk.}
\resp The multiplier and its free parameter were removed. Credits now follow the VM0033 uncertainty deduction (Eq.~2), whose only constant is the methodology's own 20\% allowance, and the registry's non-permanence buffer is applied unchanged as a separate factor. Interval width does not track 2010--2020 mangrove loss (Spearman $-0.06$ in both regions), which supports keeping the two factors separate.

\rc{The paper uses a single region and a single date.}
\resp A second region (Everglades coast) and a second Sentinel-2 date were added. A single model trained on both regions reaches cross-validated $R^2=0.37$ and 0.53; training on both dates raises the agreement of predictions across dates from 0.44 to 0.88 (Section~IV-B, Fig.~2b).

\rc{The paper has split leakage.}
\resp All models are now trained and evaluated on subgraphs induced by a single split, so no edge crosses a split boundary (Section~III-A).

\rc{The paper abandons the pre-registered Mondrian choice after seeing test results.}
\resp Both calibrations are now reported throughout. Mondrian calibration covered 87.9\% in Region~1, where only six validation blocks had enough calibration nodes, consistent with the density requirement in the spatial conformal literature. The crediting conclusion does not depend on the choice: the credited share is 0.880 (Mondrian) against 0.877 (global) in Region~1 and 0.961 against 0.974 in Region~2.

\section*{Reviewer 2}

\rc{Predictive performance remains weak; all tests yielded negative values.}
\resp Please see the response to Reviewer~1 above: after correcting the georeferencing error, test $R^2$ is 0.28 and 0.60 for the selected ensemble (Table~I).

\rc{The physics-informed validation lacks independence. The physics constraints rely on proxies, specifically NDVI and distance to tide gauges.}
\resp Both proxies were removed from the prior. The new prior uses tidal inundation frequency computed from USGS 3DEP lidar elevation and a full year of NOAA water levels, combined with a red-edge canopy index. For independent validation, 87 field cores were compared with the models and two published maps: at 25 core sites the ensemble's median error was 10.1\%, against 14.1\% and 15.1\% for the published maps and 15.0\% for a regional mean (Table~I, Section~IV-D).

\rc{Spatial leakage is present; some edges cross the training/test split.}
\resp Removed: training and evaluation now use split-induced subgraphs (Section~III-A).

\rc{The authors have not tested Random Forest or XGBoost.}
\resp Both were added, together with spatial-lag and location-only variants (Table~I). XGBoost outperformed the physics-informed GNN alone in Region~1, and their ensemble outperformed both (0.28 against 0.20 and 0.17), which indicates that the graph model contributes information the trees do not capture.

\rc{The pricing mechanism remains at the proof-of-concept stage. The multiplier parameters are arbitrary and have not been validated against market outcomes or non-permanence risks.}
\resp The multiplier was replaced by the VM0033 deduction. Against market records, the per-node deduction (median 12.7\% in Region~1) is comparable to the 11.0\% median buffer held by the seven blue carbon projects that have issued credits, of which 73.5\% have been retired; the market price reference is US\$26.96 per credit. Against non-permanence, interval width was compared with observed 2010--2020 mangrove loss. Transaction-level market outcomes are not available and are stated as a limitation.

\section*{Reviewer 3}

\rc{The introduction clearly describes the research problem \dots\ references are current and relevant \dots\ the detailed description of the methodology contributes to reproducibility.}
\resp Thank you. The revision keeps the structured description of methods and adds current references (2018--2026) on spatial validation, soil carbon controls, conformal prediction and carbon markets. All code, data-processing steps and selection rules are released to support reproducibility.

\section*{Compliance with the chairs' instructions}
The revised manuscript uses the IEEE two-column conference template, is five pages without page numbers, and embeds all fonts. Every reference is cited in the text and was verified against its DOI or source; more than half are from 2016 or later; there are no self-citations and no Wikipedia sources. Table captions are placed above the tables and figure captions below the figures, and each table and figure is cited in the text before it appears.

\end{document}
